\section{Discussion and Research Plan}
\label{sec:research-plan}

\subsection{Discussion}
\label{sec:discussion}

The main lesson of the first year concerns the benchmark rather than any one
pipeline. Every injection test on real SOSS data has to decide which light
belongs to the star, stated or not, and its verdict holds only under that
decision. For one atmosphere injected into the WASP-17~b data, this decision
reversed the ranking of the two pipelines I tested
(Section~\ref{sec:which-light}). The benchmark of
Section~\ref{sec:measured-benchmark} constrains this decision with a separate
measurement of the star. The price is a different star, field and readout. It
is designed to measure which kinds of error each method makes when the
injected light curve is known, and where in its processing they arise.

This lesson sets the order of the work. NOVA's first spectrum of the real
WASP-17~b visit is deeper than the published Ahsoka spectrum at every
wavelength, and most in the red (Section~\ref{sec:nova-real}), and its
uncertainty is not yet validated. An injection test shows how accurately each method recovers a known spectrum, which is the best guide to which real spectrum to trust, and an injector built from the WASP-17~b data alone cannot provide it. I am therefore
finishing the benchmark first. I will then use it to test NOVA and the other
pipelines and to improve NOVA, and only then refit WASP-17~b.

\subsection{Research plan}
\label{sec:plan}

The plan consists of five papers (Figure~\ref{fig:gantt}). The first contains
the benchmark, the comparison of NOVA with the other pipelines and the refit
of WASP-17~b. The second analyses five further SOSS visits as real data, the
third turns to repeated visits of the same planets, and the last two extend
NOVA to NIRSpec and to secondary eclipses.

\paragraph{Paper 1: NOVA, the benchmark and WASP-17~b.} The aim is to measure
how accurately NOVA and the extraction-first pipelines recover a known
transmission spectrum, and whether NOVA changes the conclusions about
WASP-17~b. Firstly, I will finish the benchmark of
Section~\ref{sec:measured-benchmark} and run NOVA and exoTEDRF on it, together
with Ahsoka and \texttt{transitspectroscopy} where they run end to end.
Secondly, I will improve NOVA on one set of injected atmospheres, and then
score all pipelines on a second set, held back until NOVA is fixed. Thirdly, I
will refit NOVA's spectrum of the real WASP-17~b visit, with uncertainties
from an ensemble of simulated observations, and compare it with the three
published reductions. Because the benchmark rests on a single star, I will
also propose that JWST repeat programme 4476 for several cooler stars, read
out as a SOSS time series and with more integrations at the usual position. I
aim to finish this paper before the HWO2026 conference in Paris at the end of
November 2026, where I will present this work.

\paragraph{Paper 2: Five further SOSS visits.} The aim is to find out whether
NOVA changes the atmospheric conclusions for other SOSS targets, and if so,
why. I will analyse five SOSS visits as real data. As in Paper~1, I will run
NOVA together with those pipelines of each visit's published reductions that
are publicly available, and run each of them on the benchmark as well, so that its errors on the same known case can be measured. Each visit was chosen for its own reason. WASP-39~b has
the best-studied SOSS spectrum, with six independent reductions of the same
visit, including supreme-SPOON, now exoTEDRF, and \texttt{transitspectroscopy},
and a joint analysis with three other JWST modes
\citep{FeinsteinEtAl2023,CarterEtAl2024}. Its metallicity is constrained by the water bands across the spectrum, while the fall in depth between 2 and $2.3\,\mu\mathrm{m}$ favours patchy, non-grey clouds. This is in the red range where NOVA's spectrum of WASP-17~b differs most from Ahsoka's, so the question is whether NOVA reproduces these conclusions. For WASP-96~b, published
analyses of the same visit disagree on whether the spectrum rises towards the
blue. Some find such a slope \citep{RadicaEtAl2023,TaylorEtAl2023,RadicaEtAl2026}
and favour an explanation involving small particles high in the atmosphere
\citep{TaylorEtAl2023,RadicaEtAl2026}, while another analysis does not
\citep{WangEtAl2026}. Both sides suggest differences in reduction as possible causes, and NOVA's reduction will show which of the two its spectrum supports.
The HAT-P-18~b visit contains a spot crossed by the planet during the transit
and a field star on the order-1 trace whose brightness changes, probably an
eclipsing binary \citep{FuEtAl2022}. It therefore tests how NOVA handles light
that its transit model does not describe, and Paper~3 returns to this visit
with a model of the star. HAT-P-26~b has the shallowest transit of the five,
with features of only a few hundred ppm, so it tests whether NOVA's
differences still matter for such small signals. HST and NIRSpec spectra exist
for comparison \citep{WakefordEtAl2017,GressierEtAl2025}. For WASP-80~b, NIRCam
has measured methane between 2.4 and $4\,\mu\mathrm{m}$ \citep{BellEtAl2023},
overlapping the red end of SOSS, where NOVA differs most for WASP-17~b, which
gives an independent check of NOVA's red end. I will analyse all five visits
with the same version of NOVA, fixed before the spectra are compared. Anything
learnt from them, such as a better treatment of field stars, will go into a
later version, which is again tested on the benchmark before it is used. Since
the errors measured in Paper~1 come from a hotter star, a repeat of programme
4476 on cooler stars would make this comparison firmer.

\paragraph{Paper 3: Repeated visits and stellar variability.} The aim is to
test whether repeated visits of the same planet give the same spectrum once instrumental and stellar effects are accounted for, and whether stellar effects explain features that have been attributed to planets. This matters most for
small planets around cool stars. Spots and bright regions on the star that the
planet does not cross can put false features into the spectrum
\citep{RackhamEtAl2018}, and they change, as do flares, from one visit to the
next. Comparing visits therefore tests which features persist as stellar activity changes. I will start with the two SOSS transits of GJ~9827~d
\citep{PiauletGhorayebEtAl2024}, then turn to the four of LHS~1140~b
\citep{CadieuxEtAl2024,GressierEtAl2026} and the five of TRAPPIST-1~f, each of
which contains a flare \citep{LimEtAl2026}. I will also return to HAT-P-18~b, whose spectrum rises towards the blue, to test whether spots on the star that the planet does not cross explain this rise, as \citet{FournierTondreauEtAl2024} propose, rather than a haze \citep{FuEtAl2022}. This requires building two models
into NOVA: one for flares, with their timing shared across wavelengths and
their brightness fitted at each wavelength, and one for the spots and bright
regions that the planet does not cross. I will test both models by injecting
flares and such stellar signals into the benchmark before applying them to the
real visits. Two planned MRes projects, starting in January 2027, will study
stellar variability and flares in the out-of-transit SOSS data, which would
feed directly into this paper.

\paragraph{Paper 4: Joint NIRISS and NIRSpec analysis.} The aim is to test
whether observations of the same planet with both instruments can be explained
by one transmission spectrum. NIRSpec extends the spectrum to longer
wavelengths, and where the two instruments overlap in wavelength, as with
NIRSpec's PRISM mode, their depths should agree. To test this, I will extend
NOVA and the benchmark to NIRSpec, and validate each on injected data before
comparing the real spectra. I will start with WASP-39~b and LHS~1140~b, which
have both been observed with SOSS and with NIRSpec modes that overlap it.
WASP-39~b is a well-studied gas giant with strong features \citep{CarterEtAl2024}, whereas LHS~1140~b is a small, temperate planet whose low density and lack of the strong features expected from a hydrogen-rich atmosphere favour a possible water world \citep{DamianoEtAl2024}. The benchmark needs a NIRSpec observation in which the
star is measured on its own, as programme 4476 does for SOSS. As far as I can
find, none exists yet, and I plan to propose one.

\paragraph{Paper 5: Secondary eclipses.} If time allows, I will extend NOVA
from transits to secondary eclipses, when the planet passes behind the star
and its own light briefly disappears. These signals are weak, and it would be
interesting to see whether a detector-level analysis also makes them more
reliable. The SOSS eclipse of WASP-17~b \citep{GressierEtAl2024} would be the
natural starting point.

\begin{figure*}[t]
\centering
\resizebox{\textwidth}{!}{%
\begin{ganttchart}[
  hgrid,
  vgrid,
  x unit=0.84cm,
  y unit title=0.42cm,
  y unit chart=0.36cm,
  title/.style={fill=black!6,draw=black!25},
  title label font=\bfseries\scriptsize,
  bar label font=\scriptsize,
  group label font=\bfseries\scriptsize,
  milestone label font=\scriptsize,
  bar height=0.55,
  group height=0.34,
  group peaks height=0.16,
  milestone height=0.45
]{1}{13}
\gantttitle{2026}{1}\gantttitle{2027}{4}\gantttitle{2028}{4}\gantttitle{2029}{4}\\
\gantttitle{Q4}{1}
\gantttitle{Q1}{1}\gantttitle{Q2}{1}\gantttitle{Q3}{1}\gantttitle{Q4}{1}
\gantttitle{Q1}{1}\gantttitle{Q2}{1}\gantttitle{Q3}{1}\gantttitle{Q4}{1}
\gantttitle{Q1}{1}\gantttitle{Q2}{1}\gantttitle{Q3}{1}\gantttitle{Q4}{1}\\
\ganttgroup[group/.append style={fill=ganttsoss!35,draw=ganttsoss}]
  {SOSS transmission programme}{1}{7}\\
\ganttbar[bar/.append style={fill=ganttsoss!70,draw=ganttsoss}]
  {Paper 1: NOVA, benchmark, WASP-17~b}{1}{1}\\
\ganttmilestone[milestone/.append style={fill=ganttsoss,draw=ganttsoss}]
  {HWO2026, Paris (30 Nov)}{1}\\
\ganttbar[bar/.append style={fill=ganttsoss!70,draw=ganttsoss}]
  {Paper 2: five further SOSS visits}{2}{5}\\
\ganttbar[bar/.append style={fill=ganttsoss!70,draw=ganttsoss}]
  {Paper 3: repeated visits, variability}{5}{7}\\
\ganttmilestone[milestone/.append style={fill=ganttsoss,draw=ganttsoss}]
  {SOSS analyses complete}{7}\\
\ganttgroup[group/.append style={fill=ganttgeneral!35,draw=ganttgeneral}]
  {Instrument and geometry extensions}{7}{10}\\
\ganttbar[bar/.append style={fill=ganttgeneral!70,draw=ganttgeneral}]
  {Paper 4: NIRISS and NIRSpec}{7}{9}\\
\ganttbar[bar/.append style={fill=ganttgeneral!70,draw=ganttgeneral}]
  {Paper 5: secondary eclipses}{9}{10}\\
\ganttmilestone[milestone/.append style={fill=ganttgeneral,draw=ganttgeneral}]
  {Science complete}{10}\\
\ganttgroup[group/.append style={fill=ganttwrite!35,draw=ganttwrite}]
  {Thesis}{2}{13}\\
\ganttbar[bar/.append style={fill=ganttwrite!70,draw=ganttwrite}]
  {Thesis integration}{2}{10}\\
\ganttbar[bar/.append style={fill=ganttwrite!70,draw=ganttwrite}]
  {Formal writing-up period}{11}{13}\\
\ganttmilestone[milestone/.append style={fill=ganttwrite,draw=ganttwrite}]
  {Thesis submission}{13}
\end{ganttchart}%
}
\caption{Publication and thesis plan in quarters, starting after this
assessment. The research phase ends on approximately 27 April
2029, and
the formal writing-up period runs to 28 October 2029.}
\label{fig:gantt}
\end{figure*}
